

\section*{AI agents and reproducibility}\label{ai:combined:front:start}

AI systems have recently contributed results in amplitudes~\cite{ai:Guevara:2026qzd,ai:Arkani-Hamed:2026rsz}, QCD~\cite{ai:Schwartz:2026ekw}, integrable chains~\cite{ai:Pozsgay:2026hvh} and mathematics~\cite{ai:Romera-Paredes:2024fs}. To our knowledge, language-model agents had not computed new critical exponents or CFT data; reinforcement learning has recovered known ones~\cite{ai:Kantor:2021kbx,ai:Man:2025rl}, and neural networks trained on crossing symmetry reconstruct conformal correlators and partition functions~\cite{ai:Ghosh:2026xnp,ai:Ghosh:2026efw}. This part describes how the work of this study was organized and how its results were checked.

\textit{Architecture of the collaboration.} The work was organized in layers (Fig.~\ref{ai:fig:agents}). The author set the direction: the problem, the five descriptions and key choices of method, among them computing exact asymptotic forms to guide the numerical fits and proposing the transcendental bases in which to seek exact coefficients. The author also set the rules for the clusters, approved plans and production runs, and decided between options the agents presented; the logs contain no derivation or production code by the author. OpenAI's GPT-6 astra drafted the initial high-performance codes for the clusters and computed the coefficients near the short-range end. Anthropic's Claude Code agents, in host sessions that exchanged messages and spawned subagents, reviewed and extended these codes, ran the campaigns, derived and proved closed forms, wrote the fit protocols and the research notes, and drafted the article. OpenAI's Codex, running GPT-6 astra in sessions separate from those that did the work, reviewed plans, codes and results.

\begin{figure*}[tp]
\centering
\colorlet{paperc}{brown!80!black}%
\colorlet{authc}{olive!85!black}%
\colorlet{hubc}{black!70}%
\newcommand{\agtiny}{\fontsize{6}{7}\selectfont}%
\newcommand{\agtag}[1]{\\[-0.5pt]{\ttfamily\color{black!65}#1}}%
\newcommand{\agcyc}{\tikz[baseline=-0.55ex]\draw[teal!80!black,line width=0.5pt,
  -{Stealth[length=1.8pt,width=2.2pt]}] (20:0.07) arc (20:320:0.07);}%
\newcommand{\agbadges}[3]{%
 \ifx&#2&\else\node[stk,anchor=west] at ($(#1.east)+(0.07,0.2)$) {#2};\fi
 \ifx&#3&\else\node[cdx,anchor=west] at ($(#1.east)+(0.07,-0.2)$) {\agcyc\,#3};\fi}%
\begin{tikzpicture}[x=1cm,y=1cm,font=\scriptsize,>={Stealth[length=4pt,width=3.2pt]},
 sess/.style={draw=#1,fill=#1!7,rounded corners=2pt,line width=0.6pt,
  text width=2.0cm,minimum height=1.15cm,align=center,inner sep=1.5pt},
 grp/.style={draw=#1!60,rounded corners=4pt,line width=0.5pt,fill=#1!2},
 gtitle/.style={anchor=north west,inner sep=2pt,font=\footnotesize\bfseries,text=#1},
 ext/.style={draw=#1,fill=#1!8,rounded corners=3pt,line width=0.7pt,align=center,
  inner sep=3pt,minimum height=0.95cm},
 flow/.style={->,black!75,line width=1.2pt},
 bflow/.style={<->,black!75,line width=1.2pt},
 relay/.style={->,authc,line width=0.9pt,dash pattern=on 3.5pt off 2pt},
 review/.style={<->,teal!80!black,line width=1.3pt},
 cnt/.style={font=\scriptsize,inner sep=1.2pt},
 stk/.style={draw=black!60,fill=yellow!30,rounded corners=0.8pt,inner sep=1.3pt,
  font=\agtiny\sffamily\bfseries,
  double copy shadow={shadow xshift=0.7pt,shadow yshift=0.7pt,draw=black!55,fill=yellow!12}},
 cdx/.style={draw=teal!80!black,fill=teal!10,rounded corners=0.8pt,inner sep=1.2pt,
  font=\agtiny\sffamily},
 box/.style={draw=black!40,fill=white,rounded corners=2pt,align=center,inner sep=1.5pt,
  minimum height=1.05cm},
]

\node[ext=teal!80!black,text width=3.7cm] (gpt) at (2.1,12.05)
 {{\footnotesize\textbf{GPT-6 astra}} (OpenAI)\\[-0.5pt]initial HPC codes, short-range coefficients};
\node[ext=authc,text width=3.3cm] (author) at (9.3,12.05)
 {{\footnotesize\textbf{Author}}\\[-0.5pt]directions, decisions, approvals};
\node[ext=teal!80!black,text width=5.0cm] (codex) at (14.6,12.05)
 {{\footnotesize\textbf{Codex}} (OpenAI)\\[-0.5pt]review of plans, codes and results};
\coordinate (ftop) at (0,11.05);

\draw[black!45,rounded corners=6pt,line width=0.8pt,fill=black!1] (0,2.15) rectangle (17.4,11.05);
\node[anchor=north west,inner sep=2pt,font=\footnotesize] at (0.1,11.02)
 {\textbf{Claude Code} (Anthropic): the key host sessions};

\draw[grp=cEps] (0.15,8.75) rectangle (12.25,10.6);
\node[gtitle=cEps] at (0.2,10.58) {Perturbation theory};
\draw[grp=cHT] (0.15,5.55) rectangle (12.25,8.7);
\node[gtitle=cHT] at (0.2,8.68) {Hamiltonian truncation};
\draw[grp=cNRG] (0.15,2.35) rectangle (4.0,5.4);
\node[gtitle=cNRG!80!black] at (0.2,5.38) {NRG};
\draw[grp=cMC] (4.15,2.35) rectangle (8.05,5.4);
\node[gtitle=cMC] at (4.2,5.38) {Monte Carlo};
\draw[grp=cRes] (8.2,2.35) rectangle (12.25,5.4);
\node[gtitle=cRes] at (8.25,5.38) {Resummation};
\draw[grp=black] (12.45,2.35) rectangle (17.25,10.6);
\node[gtitle=hubc] at (12.5,10.58) {Coordination and output};

\node[sess=cEps] (c81) at (1.45,9.5) {$\epsilon$ expansion, $d=1$\agtag{C81d}};
\node[sess=cEps] (c82) at (4.45,9.5) {$\epsilon$ expansion, $d=2$, general $d$\agtag{C82d}};
\node[sess=cEps] (c83) at (7.45,9.5) {$\epsilon$ expansion, $d=3$\agtag{C83d}};
\node[sess=cSR] (sri) at (10.3,9.5) {short-range end\agtag{SRI}};
\node[sess=cHT] (ht1) at (1.6,7.6) {perturbative, $d=1$\agtag{HT1dPert}};
\node[sess=cHT] (ht2) at (5.2,7.6) {perturbative, $d=2$\agtag{HT2dPert}};
\node[sess=cHT] (ht3) at (8.8,7.6) {perturbative, $d=3$\agtag{HT3dPert}};
\node[sess=cHT] (np1) at (1.6,6.2) {finite coupling, $d=1$\agtag{phi4HT1d}};
\node[sess=cHT] (np2) at (5.2,6.2) {finite coupling, $d=2$\agtag{phi4HT2d}};
\node[sess=cHT] (np3) at (8.8,6.2) {finite coupling, $d=3$\agtag{phi4HT3d}};
\node[sess=cNRG!80!black] (sb) at (1.5,3.85) {spin-boson NRG, $d=1$\agtag{spin-boson HT}};
\node[sess=cMC] (mc) at (5.5,3.85) {cluster MC, $d=1$--$3$\agtag{MC}};
\node[sess=cRes,text width=2.4cm] (res) at (9.65,3.85) {resummation, large orders, comparison of all descriptions\agtag{resum}};
\node[sess=hubc,text width=3.0cm] (sup) at (14.45,9.4) {audits, relayed decisions, coordination notes\agtag{Sup}};
\node[sess=hubc,text width=3.0cm] (pub) at (14.45,6.9) {public release: export, privacy scan, tests on two clusters\agtag{Public\_repo}};
\node[sess=paperc,text width=3.0cm] (draft) at (14.45,3.85) {writes the article\agtag{Draft}};

\agbadges{c81}{142}{74}
\agbadges{c82}{348}{39}
\agbadges{c83}{442}{34}
\agbadges{sri}{102}{11}
\agbadges{ht1}{100}{28}
\agbadges{ht2}{13}{4}
\agbadges{ht3}{3}{11}
\agbadges{np1}{1}{7}
\agbadges{np2}{}{12}
\agbadges{np3}{}{8}
\agbadges{sb}{2}{6}
\agbadges{mc}{2}{69}
\agbadges{res}{5}{}
\agbadges{sup}{24}{}
\agbadges{pub}{1}{}
\agbadges{draft}{154}{}

\draw[bflow] (12.25,9.4) -- node[cnt,above] {217} (sup.west);
\draw[bflow] (12.25,6.9) -- node[cnt,above] {470} (pub.west);
\draw[bflow] (12.25,3.85) -- node[cnt,above] {137} (draft.west);
\draw[flow] (10.75,5.55) -- ($(res.north)+(1.1,0)$);
\draw[flow] ($(mc.east)+(0.75,0)$) -- (res.west);
\draw[bflow,black!55,line width=0.8pt] (sup.south) -- node[cnt,right] {91} (pub.north);
\draw[bflow,black!55,line width=0.8pt] (pub.south) -- node[cnt,right] {146} (draft.north);
\draw[bflow,black!55,line width=0.8pt] (sup.east) -- (16.95,9.4) -- node[cnt,left,pos=0.5] {41} (16.95,3.85) -- (draft.east);

\draw[relay] (gpt.south) -- node[cnt,right,text=authc!80!black] {$\geq\AIfigGPTPackages$ code packages} (gpt.south |- ftop);
\draw[<->,authc,line width=0.9pt] (author.west) -- node[cnt,above,text=authc!80!black] {discussions} node[cnt,below,text=authc!80!black,font=\agtiny] {about half of the prompts} (gpt.east);
\draw[->,authc,line width=0.9pt] (author.south) -- node[cnt,right,text=authc!80!black] {tasks, decisions} (author.south |- ftop);
\draw[review] (codex.south) -- node[cnt,right,text=teal!60!black] {reviews and verdicts} (codex.south |- ftop);

\draw[black!45,rounded corners=4pt,fill=black!3] (0,0.05) rectangle (8.9,1.65);
\node[anchor=north west,inner sep=2pt,font=\footnotesize\bfseries] at (0.05,1.63) {Compute (Slurm)};
\node[box,text width=2.45cm] at (1.45,0.65) {\textbf{Expanse} (SDSC)};
\node[box,text width=2.45cm] at (4.25,0.65) {\textbf{symmetry} (Perimeter)};
\node[box,text width=2.45cm] at (7.15,0.65) {\textbf{two personal desktops}};
\draw[black!45,rounded corners=4pt,fill=black!3] (9.1,0.05) rectangle (17.4,1.65);
\node[anchor=north west,inner sep=2pt,font=\footnotesize\bfseries] at (9.15,1.63) {Outputs};
\node[box,text width=2.3cm] at (10.65,0.65) {research notes\agtag{\agtiny owner sessions}};
\node[box,text width=2.3cm] at (13.25,0.65) {public release\agtag{\agtiny Public\_repo}};
\node[box,text width=2.3cm] at (15.85,0.65) {the article\agtag{\agtiny Draft}};
\draw[<->,black!50,line width=1.3pt] (4.25,2.15) -- node[cnt,right,black!70] {jobs, \AIfigCoreHours} (4.25,1.65);
\draw[->,black!50,line width=1.3pt] (14.45,2.15) -- (14.45,1.65);

\begin{scope}[shift={(0.1,-0.4)}]
 \draw[bflow] (0,0) -- (0.7,0); \node[anchor=west,cnt] (l1) at (0.75,0) {messages and files, with count};
 \draw[relay] ($(l1.east)+(0.4,0)$) -- ++(0.7,0) node[anchor=west,cnt,black] (l2) {relayed by the author};
 \draw[review] ($(l2.east)+(0.4,0)$) -- ++(0.7,0) node[anchor=west,cnt,black] (l3) {Codex review};
 \node[stk,anchor=west] (k1) at ($(l3.east)+(0.4,0)$) {$n$};
 \node[anchor=west,cnt] (k1t) at ($(k1.east)+(0.06,0)$) {subagents};
 \node[cdx,anchor=west] (k2) at ($(k1t.east)+(0.35,0)$) {\agcyc\,$r$};
 \node[anchor=west,cnt] at ($(k2.east)+(0.04,0)$) {Codex rounds};
\end{scope}
\end{tikzpicture}
\caption{Collaboration of the AI agents. GPT-6 astra drafted the initial codes in discussions with the author (olive; dashed arrows: work relayed by the author). Inside the Claude Code frame, sessions that each own one topic produce the results, grouped and colored by method; Sup coordinates them, Public\_repo assembles the public release, resum combines the numbers of all descriptions, and Draft writes the article. Badges count the subagents each session spawned (yellow, \AIfigSubagents{} in all) and the rounds of review by Codex (teal, about \AIfigCodexRounds{} in all); arrows count the messages between sessions. The bottom row shows where the jobs ran and what the work produced. Of the \AIsessions{} Claude Code sessions, not counting subagents, the key ones are shown.}
\label{ai:fig:agents}
\end{figure*}


\textit{How the results were checked.} We did not rely on any single computation. Four layers of checks were used, and each result section of the article states which applied.
\begin{enumerate}
\item \emph{Review before production.} Production plans and codes were reviewed by Codex, in most cases until approved; where the author stopped a review early or a last fix was not inspected, the research notes record it.
\item \emph{Rules fixed in advance.} Reporting rules, such as error bars of four sampling standard errors, were fixed before the results where possible; where a rule was fixed afterwards and confirmed by the author, the research notes that follow say so. For the MC fits this held only in part: in $d=1$ and $d=3$ some fit choices were made or amended after data had been seen, among them the amendment under which the $d=3$ value at $s=1.6$ is accepted. Part~VI records when each rule was fixed relative to the data.
\item \emph{Independent recomputation.} Key results were recomputed by separately written codes or methods: the eight-loop coefficients in $d=2,3$, the short-range coefficients $C^{\rm SR}_5(1)$ to $C^{\rm SR}_8(1)$ and $C^{\rm SR}_2(2)$, the Hamiltonian-truncation solvers and the Wilson chain of NRG. Results reviews by Codex covered most of the final coefficients.
\item \emph{Universality.} The five descriptions compute the same $\DE$ with unrelated errors (Sec.~3.5 of the article).
\end{enumerate}
These checks caught errors. The author's own earlier calculation of the four-loop coefficient gave $C_4(1)=\CfourOneOld$; when the Claude Code agents and GPT-6 astra checked it, both found $\CfourOne$, the value in Table~1 of the article. An audit requested by the author found that an earlier version of our MC report presented a fit-model spread as a standard deviation. The checks do not cover an error shared by all codes of one description; only the comparison across descriptions does.

\clearpage
\providecommand{\href}[2]{#2}\begingroup\raggedright\begin{thebibliography}{1}

\bibitem{ai:Guevara:2026qzd}
A.~Guevara, A.~Lupsasca, D.~Skinner, A.~Strominger and K.~Weil,
  \emph{{Single-minus gluon tree amplitudes are nonzero}},
  \href{https://arxiv.org/abs/2602.12176}{{\ttfamily 2602.12176}}.

\bibitem{ai:Arkani-Hamed:2026rsz}
N.~Arkani-Hamed, F.~Calisto, N.~Ussembayev, W.W.~Zhao and Z.~Zhou,
  \emph{{Surface Water Wave Scattering and the Hydrotope}},
  \href{https://arxiv.org/abs/2606.28280}{{\ttfamily 2606.28280}}.

\bibitem{ai:Schwartz:2026ekw}
M.D.~Schwartz, \emph{{Resummation of the C-Parameter Sudakov Shoulder Using
  Effective Field Theory}},  \href{https://arxiv.org/abs/2601.02484}{{\ttfamily
  2601.02484}}.

\bibitem{ai:Pozsgay:2026hvh}
B.~Pozsgay and I.~Vona, \emph{{Bethe Ansatz with a Large Language Model}},
  \href{https://arxiv.org/abs/2603.29932}{{\ttfamily 2603.29932}}.

\bibitem{ai:Romera-Paredes:2024fs}
B.~Romera-Paredes, M.~Barekatain, A.~Novikov, M.~Balog, M.P.~Kumar, E.~Dupont
  et~al., \emph{{Mathematical discoveries from program search with large
  language models}},
  \href{https://doi.org/10.1038/s41586-023-06924-6}{\emph{Nature} {\bfseries
  625} (2024) 468}.

\bibitem{ai:Kantor:2021kbx}
G.~K{\'a}ntor, V.~Niarchos and C.~Papageorgakis, \emph{{Solving Conformal Field
  Theories with Artificial Intelligence}},
  \href{https://doi.org/10.1103/PhysRevLett.128.041601}{\emph{Phys. Rev. Lett.}
  {\bfseries 128} (2022) 041601}
  [\href{https://arxiv.org/abs/2108.08859}{{\ttfamily 2108.08859}}].

\bibitem{ai:Man:2025rl}
H.~Man, C.~Wang, J.-R.~Li, Y.~Tian and S.-G.~Chen, \emph{{Autonomous discovery
  of the Ising model's critical parameters with reinforcement learning}},
  \href{https://doi.org/10.1088/1742-5468/ae22ea}{\emph{J. Stat. Mech.}
  {\bfseries 2025} (2025) 123410}
  [\href{https://arxiv.org/abs/2601.05577}{{\ttfamily 2601.05577}}].

\bibitem{ai:Ghosh:2026xnp}
K.~Ghosh, S.~Kumar, V.~Niarchos and A.~Stergiou, \emph{{Neural spectral bias
  and conformal correlators I introduction and applications}},
  \href{https://doi.org/10.1016/j.nuclphysb.2026.117609}{\emph{Nucl. Phys. B}
  {\bfseries 1030} (2026) 117609}
  [\href{https://arxiv.org/abs/2604.18686}{{\ttfamily 2604.18686}}].

\bibitem{ai:Ghosh:2026efw}
K.~Ghosh, S.~Kumar, V.~Niarchos and A.~Stergiou, \emph{{Neural Spectral Bias
  and Conformal Correlators II: Modular and Annulus Bootstrap}},
  \href{https://arxiv.org/abs/2607.24913}{{\ttfamily 2607.24913}}.

\end{thebibliography}\endgroup
